\documentclass[12pt]{article}
\usepackage{amsmath, amssymb}
\usepackage[margin=1in]{geometry}
\setlength{\parskip}{6pt}
\title{QUBO Formulation: Single-Track Railway Dispatching Problem}
\date{}
\begin{document}
\maketitle

\subsection*{Single-Track Railway Dispatching Problem}

\paragraph{Problem Statement.}
Given a set of $N$ trains $\mathcal{T} = \{0, \dots, N-1\}$ that must traverse a shared single-track segment, each train $t \in \mathcal{T}$ is associated with a set of $M_t$ feasible timing options $\mathcal{O}_t = \{0, \dots, M_t-1\}$. Each option $o \in \mathcal{O}_t$ incurs a non-negative delay cost $c_{t,o}$. Certain pairs of options from different trains are mutually exclusive due to track occupancy conflicts, denoted by the conflict set $\mathcal{C} \subseteq \mathcal{T} \times \mathcal{O} \times \mathcal{T} \times \mathcal{O}$. The task is to select exactly one timing option for every train such that no two selected options conflict, thereby minimizing the total dispatch cost.

\paragraph{Decision Variables.}
For each train $t \in \mathcal{T}$ and each feasible option $o \in \mathcal{O}_t$, introduce a binary decision variable $x_{t,o} \in \{0,1\}$. The variable $x_{t,o}$ equals $1$ if train $t$ is assigned timing option $o$, and $0$ otherwise. The full decision vector is $\mathbf{x} = (x_{t,o})_{t \in \mathcal{T}, o \in \mathcal{O}_t} \in \{0,1\}^n$, where $n = \sum_{t \in \mathcal{T}} M_t$.

\paragraph{QUBO Derivation.}
\textbf{Step 1 --- Original Objective.}
The problem requires minimizing the sum of selected delay costs plus a large penalty $P > 0$ for each selected conflicting pair:
\begin{align}
\min_{\mathbf{x}} \; f(\mathbf{x}) &= \sum_{t \in \mathcal{T}} \sum_{o \in \mathcal{O}_t} c_{t,o} x_{t,o} + P \sum_{((t_1,o_1),(t_2,o_2)) \in \mathcal{C}} x_{t_1,o_1} x_{t_2,o_2}.
\end{align}

\textbf{Step 2 --- Constraints.}
Each train must be assigned exactly one timing option:
\begin{align}
\sum_{o \in \mathcal{O}_t} x_{t,o} = 1, \quad \forall t \in \mathcal{T}.
\end{align}

\textbf{Step 3 --- Penalty Term Expansion.}
We enforce the constraint by adding a quadratic penalty with weight $\lambda > 0$:
\begin{align}
\mathcal{P}(\mathbf{x}) &= \lambda \sum_{t \in \mathcal{T}} \left( \sum_{o \in \mathcal{O}_t} x_{t,o} - 1 \right)^2.
\end{align}
Expanding the square and applying the idempotent property $x_{t,o}^2 = x_{t,o}$ for binary variables yields:
\begin{align}
\left( \sum_{o \in \mathcal{O}_t} x_{t,o} - 1 \right)^2 &= \sum_{o \in \mathcal{O}_t} x_{t,o}^2 + \sum_{o \in \mathcal{O}_t} \sum_{\substack{o' \in \mathcal{O}_t \\ o' \neq o}} x_{t,o} x_{t,o'} - 2 \sum_{o \in \mathcal{O}_t} x_{t,o} + 1 \nonumber \\
&= \sum_{o \in \mathcal{O}_t} x_{t,o} + \sum_{o \neq o'} x_{t,o} x_{t,o'} - 2 \sum_{o \in \mathcal{O}_t} x_{t,o} + 1 \nonumber \\
&= -\sum_{o \in \mathcal{O}_t} x_{t,o} + \sum_{o \neq o'} x_{t,o} x_{t,o'} + 1.
\end{align}
Summing over all trains, the total penalty contribution to the Hamiltonian is:
\begin{align}
\mathcal{P}(\mathbf{x}) = \lambda \sum_{t \in \mathcal{T}} \left( -\sum_{o \in \mathcal{O}_t} x_{t,o} + \sum_{o \neq o'} x_{t,o} x_{t,o'} + 1 \right).
\end{align}

\textbf{Step 4 --- Combined QUBO Hamiltonian.}
Adding the original objective and the penalty term gives the unconstrained QUBO Hamiltonian $H(\mathbf{x})$:
\begin{align}
H(\mathbf{x}) &= \sum_{t \in \mathcal{T}} \sum_{o \in \mathcal{O}_t} c_{t,o} x_{t,o} + P \sum_{((t_1,o_1),(t_2,o_2)) \in \mathcal{C}} x_{t_1,o_1} x_{t_2,o_2} \nonumber \\
&\quad + \lambda \sum_{t \in \mathcal{T}} \left( -\sum_{o \in \mathcal{O}_t} x_{t,o} + \sum_{o \neq o'} x_{t,o} x_{t,o'} + 1 \right).
\end{align}

\textbf{Step 5 --- Q Matrix Entries and Offset.}
Collecting linear and quadratic terms, the Hamiltonian takes the standard form $H(\mathbf{x}) = \sum_{i \le j} Q_{i,j} x_i x_j + c$. Mapping each pair $(t,o)$ to a unique index $i \in \{0,\dots,n-1\}$, the matrix entries and constant offset are:
\begin{align}
Q_{(t,o),(t',o')} &= \begin{cases} 
c_{t,o} - \lambda & \text{if } t=t' \text{ and } o=o', \\
P/2 & \text{if } ((t,o),(t',o')) \in \mathcal{C} \text{ and } (t,o) \neq (t',o'), \\
\lambda & \text{if } t=t' \text{ and } o \neq o', \\
0 & \text{otherwise},
\end{cases} \\
c &= \lambda N.
\end{align}

\paragraph{Penalty Weight Justification.}
The penalty weight $\lambda$ must be chosen sufficiently large to ensure that any violation of the exactly-one constraint is strictly more costly than the maximum possible reduction in the original objective. The maximum possible decrease in delay cost from selecting a suboptimal option is bounded by $\max_{t,o} c_{t,o}$. Therefore, selecting $\lambda > \max_{t,o} c_{t,o}$ guarantees that any infeasible assignment (where $\sum_{o} x_{t,o} \neq 1$ for some $t$) incurs a penalty of at least $\lambda$, which dominates any feasible objective value.

\paragraph{Correctness Argument.}
Minimizing $H(\mathbf{x})$ over $\{0,1\}^n$ recovers the optimal dispatch schedule because (i) for any feasible assignment, all constraint terms vanish, reducing $H(\mathbf{x})$ to the original delay cost plus conflict penalties; (ii) for any infeasible assignment, at least one train is assigned zero or multiple options, incurring a penalty of at least $\lambda$. Since $\lambda$ exceeds the maximum possible objective value of any feasible solution, every infeasible point has strictly higher energy than the best feasible point. Consequently, the global minimum of $H(\mathbf{x})$ must be feasible and, by construction, minimizes the original objective function.

\paragraph{Illustrative Example.}
Consider a concrete instance with $N=2$ trains and $M_0=M_1=2$ options per train. The variables are mapped to indices $0 \to (0,0)$, $1 \to (0,1)$, $2 \to (1,0)$, $3 \to (1,1)$. The conflict set is $\mathcal{C} = \{((0,0),(1,0)), ((0,1),(1,1))\}$. Using the general formulas above, the concrete objective terms are $c_{0,0}x_0 + c_{0,1}x_1 + c_{1,0}x_2 + c_{1,1}x_3 + P x_0 x_2 + P x_1 x_3$. The penalty terms for the exactly-one constraints expand to $-\lambda x_0 - \lambda x_1 + 2\lambda x_0 x_1 - \lambda x_2 - \lambda x_3 + 2\lambda x_2 x_3$, with an offset of $2\lambda$. The resulting Hamiltonian is:
\begin{align}
H(\mathbf{x}) &= (c_{0,0}-\lambda)x_0 + (c_{0,1}-\lambda)x_1 + (c_{1,0}-\lambda)x_2 + (c_{1,1}-\lambda)x_3 \nonumber \\
&\quad + P x_0 x_2 + P x_1 x_3 + 2\lambda x_0 x_1 + 2\lambda x_2 x_3 + 2\lambda.
\end{align}
For a specific numerical instantiation, let $c_{0,0}=1, c_{0,1}=5, c_{1,0}=3, c_{1,1}=2, P=100, \lambda=10$. The feasible assignments are $(x_0,x_1,x_2,x_3) = (1,0,0,1)$ with energy $1+2 + 2(10) = 23$, and $(0,1,1,0)$ with energy $5+3 + 2(10) = 28$. The global minimum is $\mathbf{x}^* = (1,0,0,1)$, corresponding to selecting option $0$ for train $0$ and option $1$ for train $1$, yielding a minimum feasible delay cost of $3$.

\end{document}